% Provenance: COMPATIBILIDAD_HOLONOMIA_E_INFORMACION.md, sections 2--9.
% Scientific development preserved; editorial records and correspondence are separate.
\section{Constitutive compatibility, relative holonomy and operational information}
\label{md:comp:capitulo}

\subsection{Differential reciprocity of speed and impedance}
\label{md:comp:reciprocidad}

The constitutive antecedent in Section~\ref{iii:sec:constitutiva}
constructs, on the retained sheets of its representation,
\begin{equation}
 \begin{gathered}
 T=e^{-xI-y\mathcal R},\qquad
 r_\pm=f\bigl(e^{-30(x\pm y)}\bigr),\\
 V=(I-T^{90})(I-T^{120})^{-1}=r_+P_++r_-P_-.
 \end{gathered}
 \label{md:comp:constitutiva}
\end{equation}
where
\[
 f(s)=\frac{1-s^3}{1-s^4},\qquad x>|y|.
\]
The sheet projectors remain fixed. The marked physical chamber
\(x>y>0\) and its conjugate lie within this analytic collection.
Define the normalized logarithmic coordinates
\begin{equation}
 c_\ell=\log\widehat c=-\log r_+-\log r_-,\qquad
 z_\ell=\log\widehat Z=\log r_--\log r_+.
 \label{md:comp:logaritmos}
\end{equation}
These letters denote neither new constants nor the counterangle \(C^*\).
Let
\[
 u=30(x+y),\qquad v=30(x-y),\qquad
 g(w)=\log f(e^{-w}),\qquad a=g'(u),\quad b=g'(v).
\]

\begin{proposition}[Differential constitutive reciprocity]
\label{md:comp:prop-reciprocidad}
In the domain \(x>|y|\), one has
\begin{equation}
 g'(w)=\frac{3}{e^{3w}-1}-\frac{4}{e^{4w}-1}>0,
 \qquad
 D(c_\ell,z_\ell)=-30
 \begin{pmatrix}a+b&a-b\\a-b&a+b\end{pmatrix}.
 \label{md:comp:jacobiano}
\end{equation}
The eigenvalues of the Jacobian are \(-60a\) and \(-60b\).
It is therefore symmetric and negative definite, its determinant is
\(3600ab>0\), and
\begin{equation}
 \boxed{\partial_y\log\widehat c=\partial_x\log\widehat Z.}
 \label{md:comp:reciprocidad-cruzada}
\end{equation}
\end{proposition}
\begin{proof}
Differentiate the identities
\(c_\ell=-g(u)-g(v)\) and \(z_\ell=g(v)-g(u)\).
Positivity of \(g'\) follows from the strict decrease, in \(k>0\),
of \(k/(e^{kw}-1)\). It also follows from the positive rational
expression in Section~\ref{md:comp:restitucion}. The resulting matrix
has the stated eigenvalues and determinant; equality of its off-diagonal
entries gives \eqref{md:comp:reciprocidad-cruzada}.
\end{proof}

The equality relates the response of speed to oriented separation
with the response of impedance to the common coordinate.
Both sensitivities arise from the same spectral function. This is a
reciprocity of the present constitutive collection; it is not identified
without proof with a law of thermodynamics or macroscopic electromagnetism.

\subsection{Spectral potential and the depth-two connection with Catalan}
\label{md:comp:potencial}

\subsubsection{Construction of the potential}
\label{md:comp:construccion-potencial}

From the already constructed channels, define
\begin{equation}
 \mathscr F(x,y)=\sum_{\sigma=\pm1}\sum_{n\geq1}
 \left[
 \frac{e^{-120n(x+\sigma y)}}{120n^2}
 -\frac{e^{-90n(x+\sigma y)}}{90n^2}
 \right].
 \label{md:comp:serie-potencial}
\end{equation}

\begin{proposition}[Potential of the two logarithmic responses]
\label{md:comp:prop-potencial}
The series \eqref{md:comp:serie-potencial} defines on \(x>|y|\) a
potential satisfying
\begin{equation}
 \boxed{\partial_x\mathscr F=c_\ell,\qquad
        \partial_y\mathscr F=z_\ell.}
 \label{md:comp:gradiente-potencial}
\end{equation}
Its Hessian is the Jacobian in
\eqref{md:comp:jacobiano}; consequently, \(\mathscr F\) is
strictly concave. The additive constant is fixed by
\(\mathscr F\longrightarrow0\) as \(x-|y|\longrightarrow\infty\).
\end{proposition}
\begin{proof}
On every compact subset of the domain \(x>|y|\), the quantities
\(x\pm y\) have a positive minimum. A power of \(n\) multiplied by
a decreasing exponential is summable; hence the series and its
derivatives of every fixed order converge uniformly on that compact set.
Termwise differentiation, using
\(-\log(1-z)=\sum_{n\geq1}z^n/n\), yields the two identities in
\eqref{md:comp:gradiente-potencial}. The Hessian then coincides with
\eqref{md:comp:jacobiano}, which is negative definite.
The normalization at infinity is that of the displayed series.
\end{proof}

\(\mathscr F\) is a mathematical response potential. No energy units
are assigned to it, nor is it declared to be a thermodynamic potential.

\subsubsection{Realization through a common spectral moment}
\label{md:comp:momento-espectral}

Let
\[
 D_2(z)=\sum_{n\geq1}\frac{z^n}{n^2}.
\]
Its recognition as the dilogarithm follows this series definition.
For \(|z|<1\), it satisfies
\((z\partial_z)^2D_2(z)=z/(1-z)\). In this notation,
\begin{equation}
 \mathscr F=
 \sum_{\sigma\in\{+,-\}}
 \left[\frac{D_2(q_\sigma^{120})}{120}
       -\frac{D_2(q_\sigma^{90})}{90}\right].
 \label{md:comp:potencial-dos-momentos}
\end{equation}
The oriented Catalan moment constructed in
Section~\ref{iii:sec:catalan} satisfies
\begin{equation}
 \mathcal C_2(s)=\operatorname{Im}D_2(is)
 =\sum_{k\geq0}\frac{(-1)^ks^{2k+1}}{(2k+1)^2},\qquad
 \mathcal C_2(1)=G_{\rm Cat}.
 \label{md:comp:catalan-momento}
\end{equation}

\begin{proof}
Even powers of \(i\) are real, and odd powers have imaginary part
\((-1)^k\). This gives the series in
\eqref{md:comp:catalan-momento}, which converges absolutely even at
\(s=1\). For \(0\leq s<1\), applying \(s\partial_s\) twice termwise
returns \(s/(1+s^2)\), the oriented coefficient of the quarter-turn
resolvent already present in the corpus.
If used at \(s=1\), the differentiated identity is interpreted through
analytic continuation or an Abel limit: the twice-differentiated series
does not converge in the ordinary sense at that endpoint.
The identification \(\mathcal C_2(1)=G_{\rm Cat}\) follows directly
from the original absolutely convergent series.
\end{proof}

Thus the Catalan moment and the constitutive potential are brought
together through a common depth-two spectral moment, evaluated on
different supports: the oriented quarter-turn and the contractions of
the \(90/120\) sectors. This composition neither identifies Catalan
with \(\mathscr F\) nor asserts that its single endpoint value determines
all channels. The analytic collection and its arguments retain information
that an isolated value does not contain.

\subsubsection{Compatibility with amplification}
\label{md:comp:amplificacion-potencial}

For \(T_m=T\otimes I_{9^m}\) and
\(\tau_m=\operatorname{Tr}/(2\cdot9^m)\), one has exactly
\begin{equation}
 \mathscr F_m=2\tau_m\left[
 \frac{D_2(T_m^{120})}{120}
 -\frac{D_2(T_m^{90})}{90}\right]=\mathscr F.
 \label{md:comp:potencial-amplificado}
\end{equation}
The multiplicity \(9^m\) cancels the normalization. The same argument
preserves the derivatives. This establishes amplification of the
representation, not every conceivable refinement of the TPK.
The temporal operator is not replaced by a list of eigenvalues lacking
sheet markings.

For every smooth closed loop in this analytic collection,
\begin{equation}
 \oint(c_\ell\,dx+z_\ell\,dy)=0.
 \label{md:comp:integral-cerrada}
\end{equation}
This equality concerns the two-coordinate constitutive projection.
It does not imply that the enriched history or its operator holonomy
is trivial.

\subsection{Exact reconstruction and finite sensitivity}
\label{md:comp:restitucion}

\begin{proposition}[Analytic reconstruction of the angular coordinates]
\label{md:comp:prop-restitucion}
With \(L=\log(16/9)\), the map
\((x,y)\longmapsto(c_\ell,z_\ell)\) is a real-analytic diffeomorphism
from the domain \(x>|y|\) onto
\begin{equation}
 0<c_\ell+z_\ell<L,\qquad 0<c_\ell-z_\ell<L.
 \label{md:comp:imagen-logaritmica}
\end{equation}
\end{proposition}
\begin{proof}
The function \(g\) is strictly increasing from \(\log(3/4)\) to \(0\).
The combinations
\[
 c_\ell+z_\ell=-2g(u),\qquad c_\ell-z_\ell=-2g(v)
\]
recover \(u,v\) through \(g^{-1}\). Then
\(x=(u+v)/60\) and \(y=(u-v)/60\).
Injectivity is preserved under restriction to the image actually
produced by HMT; no assertion is made that HMT produces the entire cone.
\end{proof}

With \(s=e^{-w}\), the derivative has the rational form
\begin{equation}
 a(w)=g'(w)=
 \frac{s^3(s^2+2s+3)}{(1+s+s^2)(1+s+s^2+s^3)},
 \qquad 0<a(w)<\frac12.
 \label{md:comp:derivada-racional}
\end{equation}
The function is strictly decreasing, with \(a(0+)=1/2\), and
\(a(w)\sim3e^{-3w}\) at infinity. To prove the decrease, use
\begin{equation}
 a'(w)=-\frac{9}{4\sinh^2(3w/2)}
        +\frac{16}{4\sinh^2(2w)}<0,
 \label{md:comp:decrecimiento}
\end{equation}
since \(k/\sinh(kw/2)\) is strictly decreasing in \(k>0\).
Moreover,
\begin{equation}
 \frac12e^{-3w}<a(w)<3e^{-3w}.
 \label{md:comp:cotas-exponenciales}
\end{equation}
The lower bound follows by multiplying the denominators and using
\begin{equation}
 \begin{split}
 &2(s^2+2s+3)-(1+s+s^2)(1+s+s^2+s^3)\\
 &\hspace{1cm}=(1-s)(5+7s+6s^2+3s^3+s^4)>0.
 \end{split}
 \label{md:comp:polinomio-cota}
\end{equation}
The upper bound follows from
\[
 3(1+s+s^2)(1+s+s^2+s^3)-(s^2+2s+3)>0.
\]

If \(m=30(x-|y|)\) and \(M=30(x+|y|)\), the norm of the inverse
Jacobian is \(1/[60a(M)]\). To distinguish the vector-valued map from
the scalar potential, write \(\mathbf F=(c_\ell,z_\ell)\).

\begin{proposition}[Finite reconstruction bounds]
\label{md:comp:prop-sensibilidad}
On a convex region in which both arguments \(u,v\) belong to
\([m_0,M_0]\subset(0,\infty)\), the following bounds hold:
\begin{equation}
 60a(M_0)\|p-q\|
 \leq\|\mathbf F(p)-\mathbf F(q)\|
 \leq60a(m_0)\|p-q\|.
 \label{md:comp:cotas-finitas}
\end{equation}
\end{proposition}
\begin{proof}
Integrate the Hessian of \(\mathscr F\), namely the Jacobian of
\(\mathbf F\), along the segment. For the lower bound, use
\[
 -\langle\mathbf F(p)-\mathbf F(q),p-q\rangle
 \geq60a(M_0)\|p-q\|^2
\]
and the Cauchy--Schwarz inequality. The upper bound uses the maximum
norm of the Jacobian. These are finite-error bounds on the specified
domain.
\end{proof}

At \(y=0\), as \(x\longrightarrow\infty\), both directions contract
equally: the matrix condition number is one, but the norm of the inverse
grows as \(e^{90x}/180\). Exact invertibility and recovery robust to
absolute error are therefore different properties. As
\(|y|\longrightarrow x\) with fixed \(x>0\), there is no differential
singularity: one derivative tends to \(1/2\), and the other to
\(a(60x)>0\).

\subsection{The quarter-turn arises from the dodecaphasic cycle}
\label{md:comp:cuarto-giro}

The construction of \eqref{iii:eq:quarter-intertwiner} is recovered.
In \(\mathbb R^{12}\), let
\[
 C_{12}e_j=e_{j+1\bmod12},\qquad
 \Pi_4=P_4^{\rm cyc}\frac{I-C_{12}^6}{2},
\]
and
\begin{equation}
 \begin{gathered}
 u_0=(1,0,-1,0,1,0,-1,0,1,0,-1,0)^{\mathsf T},\qquad
 v_0=C_{12}u_0,\\
 W=(u_0,v_0)/\sqrt6.
 \end{gathered}
 \label{md:comp:isometria-ciclica}
\end{equation}
The symbol \(\Pi_4\) distinguishes the projector from the action
coordinates \(Q,P\). The following identities hold:
\begin{equation}
 \begin{gathered}
 W^{\mathsf T}W=I,\qquad WW^{\mathsf T}=\Pi_4,\qquad
 C_{12}W=WJ_0,\\
 J_0=\begin{pmatrix}0&-1\\1&0\end{pmatrix},\qquad J_0^2=-I.
 \end{gathered}
 \label{md:comp:entrelazador-ciclico}
\end{equation}
In particular, \(C_{12}^3W=-WJ_0\). The plane thus has an explicit
intertwiner with the cycle; no external rotation has been introduced to
fit a response.

The radial antecedent of \eqref{rec:eq:entrelazamiento} realizes
\[
 S_\zeta=\operatorname{diag}(e^{\zeta/2},e^{-\zeta/2}),
 \qquad J_\zeta=S_\zeta J_0S_\zeta^{-1}.
\]
On this plane, define
\begin{equation}
 U_\zeta(\theta)=S_\zeta e^{-\theta J_0}S_\zeta^{-1}
 =\begin{pmatrix}
 \cos\theta&e^\zeta\sin\theta\\
 -e^{-\zeta}\sin\theta&\cos\theta
 \end{pmatrix}.
 \label{md:comp:transporte-plano}
\end{equation}
The sign of \(\theta\) fixes the chosen orientation. The generated
ratio \(\xi=C^*/A\), with \(|\xi|<1\) and both angles expressed in the
same unit, satisfies \(\zeta=\operatorname{artanh}\xi\).
In the conjugate orientation,
\((\varepsilon,\xi)\mapsto(-\varepsilon,-\xi)\) preserves the positive
action section through the oriented reconstruction map of the corpus;
\(\zeta\) changes sign.

\subsection{Relative loop of two transports}
\label{md:comp:lazo}

Let
\begin{equation}
 U=U_{\zeta_s}(\theta),\qquad
 V=U_{\zeta_r}(\vartheta),\qquad
 H_{\rm lazo}=UVU^{-1}V^{-1}.
 \label{md:comp:def-lazo}
\end{equation}
Products act from right to left.
Here \(\vartheta\) denotes the second phase, to distinguish it from
the regional coordinate \(\varphi\). Put
\[
 \delta=\zeta_s-\zeta_r,\qquad
 \chi=\sinh\delta\,\sin\theta\,\sin\vartheta.
\]

\begin{proposition}[Trace and spectrum of the relative loop]
\label{md:comp:prop-lazo}
The composition \eqref{md:comp:def-lazo} satisfies
\begin{equation}
 \boxed{\begin{gathered}
 \det H_{\rm lazo}=1,\qquad
 \operatorname{tr}H_{\rm lazo}=2+4\chi^2,\\
 H_{\rm lazo}=I\ \Longleftrightarrow\ \chi=0.
 \end{gathered}}
 \label{md:comp:traza-lazo}
\end{equation}
Its eigenvalues are
\begin{equation}
 \lambda_\pm=(\sqrt{1+\chi^2}\pm|\chi|)^2
 =\exp\bigl(\pm2\operatorname{arsinh}|\chi|\bigr).
 \label{md:comp:espectro-lazo}
\end{equation}
\end{proposition}
\begin{proof}
Under simultaneous conjugation by \(S_{\zeta_r}^{-1}\), the operator
\(V\) becomes \(R(-\vartheta)\), and \(U\) has terms
\(e^{\pm\delta}\sin\theta\). Direct multiplication gives, in that
reference,
\[
 UV-VU=-2\chi\operatorname{diag}(1,-1).
\]
Since
\[
 H_{\rm lazo}-I=(UV-VU)U^{-1}V^{-1},
\]
the equivalence with \(\chi=0\) follows. Multiplication, or the
trace identity for unimodular matrices, gives
\(\operatorname{tr}H_{\rm lazo}=2+4\chi^2\).
The characteristic polynomial is
\(\lambda^2-(2+4\chi^2)\lambda+1\), whose roots are those in
\eqref{md:comp:espectro-lazo}.
\end{proof}

This is a relative holonomy of the declared composition.
It is not identified with \(\Gamma_9\) or with a complete TPK update
on the basis of similar terminology.

\subsubsection{Specialization fully determined by the conjugate pair}
\label{md:comp:lazo-conjugado}

Let \(U_+=U_\zeta(\pi/2)\) and
\(U_-=U_{-\zeta}(\pi/2)\). Both satisfy
\(U_+^2=U_-^2=-I\), but
\begin{equation}
 \begin{aligned}
 U_+U_-&=-\operatorname{diag}(e^{2\zeta},e^{-2\zeta}),\\
 H_{\rm lazo}=(U_+U_-)^2
 &=\operatorname{diag}(e^{4\zeta},e^{-4\zeta}).
 \end{aligned}
 \label{md:comp:lazo-diagonal}
\end{equation}
Writing \(\xi=\tanh\zeta=C^*/A\), with both angles expressed in the
same unit, gives
\begin{equation}
 \begin{aligned}
 H_{\rm lazo}
 &=\operatorname{diag}\left(
 \left[\frac{1+\xi}{1-\xi}\right]^2,
 \left[\frac{1-\xi}{1+\xi}\right]^2\right),\\
 \operatorname{tr}H_{\rm lazo}-2
 &=\frac{16\xi^2}{(1-\xi^2)^2}.
 \end{aligned}
 \label{md:comp:lazo-razon-angular}
\end{equation}
The formula for \(H_{\rm lazo}\) composes the quarter-turn and the
conjugate shapes of the same pair. Its expression is exact in the
generated angular ratio \(\xi\).

On the positive action branch of the corpus,
\(A_{\deg}=1000\alpha\) and
\(C^*_{\deg}=2(\eta_{\rm ret}+\alpha)\), with
\(\hbar_{\rm ret}=s_0\eta_{\rm ret}\). Hence,
\begin{equation}
 \begin{aligned}
 \xi&=\frac{\eta_{\rm ret}+\alpha}{500\alpha},\\
 \rho_{\rm lazo}=e^{4\zeta}
 &=\left(\frac{501\alpha+\eta_{\rm ret}}
                 {499\alpha-\eta_{\rm ret}}\right)^2,
 \qquad 0<\eta_{\rm ret}<499\alpha.
 \end{aligned}
 \label{md:comp:multiplicador-lazo}
\end{equation}
The coefficients \(501\) and \(499\) are \(500\pm1\);
\(500\) comes from the factor \(1000/2\) in the angular pair and is
not a new parameter. With the oriented axes retained, reconstruction is
\begin{equation}
 \boxed{
 \eta_{\rm ret}=\alpha\left[
 500\frac{\sqrt{\rho_{\rm lazo}}-1}
          {\sqrt{\rho_{\rm lazo}}+1}-1\right].}
 \label{md:comp:restitucion-accion}
\end{equation}
The derivative
\[
 \frac{d\eta_{\rm ret}}{d\rho_{\rm lazo}}
 =\frac{500\alpha}
 {\sqrt{\rho_{\rm lazo}}(\sqrt{\rho_{\rm lazo}}+1)^2}
\]
is positive. The chamber \(\eta_{\rm ret}>0\) corresponds to
\(\rho_{\rm lazo}>(501/499)^2\); the boundary
\(\eta_{\rm ret}\longrightarrow499\alpha\) corresponds to
\(\rho_{\rm lazo}\longrightarrow\infty\).
This is an inverse readout of the dimensionless action coefficient
from relative transport and the already constructed coordinate
\(\alpha\). It retains the scale \(s_0\) as an antecedent: a
unimodular shape transformation does not by itself determine that
dimensional scale. The conjugate orientation inverts \(\rho_{\rm lazo}\)
and transports the action label; it does not turn
\(\eta_{\rm ret}\) into negative action.

\subsubsection{Loop iteration and oriented accumulation}
\label{md:comp:iteracion}

Let
\[
 \rho_{\rm lazo}=e^{4\zeta}
 =\left(\frac{1+\xi}{1-\xi}\right)^2.
\]
For every integer \(n\),
\begin{equation}
 \begin{gathered}
 H_{\rm lazo}^n
 =\operatorname{diag}(\rho_{\rm lazo}^{n},\rho_{\rm lazo}^{-n}),
 \\
 Q_n=\rho_{\rm lazo}^{n}Q_0,\quad
 P_n=\rho_{\rm lazo}^{-n}P_0.
 \end{gathered}
 \label{md:comp:iteracion-diagonal}
\end{equation}
The proof is multiplication of diagonal matrices, extended to \(n<0\)
through the inverse. The product \(Q_nP_n=Q_0P_0\) and symplectic
area are preserved, whereas the state changes for \(\zeta\ne0\) and
\(X_0\ne0\). In the four-operation protocol, the declared phase
increments sum to zero at each repetition; that phase accounting does
not determine the total state transport.

For \(\zeta\ne0\) and \(Q_0\ne0\), one recovers exactly
\[
 n=\frac{\log|Q_n/Q_0|}{4\zeta}.
\]
If \(Q_0=0\), the condition \(P_0\ne0\) allows use of
\[
 n=-\frac{\log|P_n/P_0|}{4\zeta}.
\]
The readout requires knowledge of the preparation and preservation of a
reference. The multiplicative extension
\(H_{\rm lazo}^{n+m}=H_{\rm lazo}^{n}H_{\rm lazo}^{m}\)
corresponds to an additive logarithmic increment \(4(n+m)\zeta\).
Reversing the loop changes the sign of that increment.

Thus a fixed repeated protocol admits a response that records its
iteration. This is a mathematical realization of operational accumulation
in the composition studied. The integer \(n\) is not identified with
all carry or memory of \(\Gamma_9\), nor is zero entropy for the complete
dynamics inferred from periodicity of its phase protocol. If the initial
amplitude is unknown and only the final state is available, \(n\) is
not determined by that isolated readout.

\subsubsection{Explicit lift to the twelve phases}
\label{md:comp:levantamiento}

Let
\[
 \widetilde S_\zeta=I-\Pi_4+WS_\zeta W^{\mathsf T},\qquad
 L_\pm=\widetilde S_{\pm\zeta}C_{12}^3
       \widetilde S_{\pm\zeta}^{-1}.
\]
On \(\operatorname{Ran}\Pi_4\), these act as \(U_+,U_-\),
respectively; on the complement, both act as \(C_{12}^3\). Hence,
\begin{equation}
 L_+L_-L_+^{-1}L_-^{-1}
 =I-\Pi_4+WH_{\rm lazo}W^{\mathsf T}.
 \label{md:comp:levantamiento-lazo}
\end{equation}
The complementary action cancels. The trace of the twelve-phase
operator is \(10+\operatorname{tr}H_{\rm lazo}\).
The intertwining and composition are exact and have been checked
rationally using the integer vectors \(u_0,v_0\).
This lift brings together the dodecaphasic antecedent and elliptic
transport; it does not assert that its sequence is in itself the
original autonomous TPK update.

\subsubsection{Refinement and limit}
\label{md:comp:limite}

The amplification
\(H_{{\rm lazo},m}=H_{\rm lazo}\otimes I_{9^m}\) preserves
\[
 \frac{1}{9^m}\operatorname{Tr}H_{{\rm lazo},m}
 =\operatorname{Tr}H_{\rm lazo},
\]
as well as the norm and spectrum, apart from multiplicities.
With the inclusions \(v\mapsto v\otimes e_0\), one has
\[
 H_{{\rm lazo},m+1}\iota_m=\iota_m H_{{\rm lazo},m}.
\]
The uniform operator defines a bounded extension on the Hilbert-space
inductive limit, with the same norm. This is the limit of this
amplification; it does not replace any other refinement system in
the corpus.

\subsection{Distinguishing a chart change from a relative effect}
\label{md:comp:carta-efecto}

The coordinate transport
\[
 T_n=S_{\zeta_{n+1}}R(-\theta_n)S_{\zeta_n}^{-1}
\]
telescopes:
\begin{equation}
 T_{N-1}\cdots T_0
 =S_{\zeta_N}R\left(-\sum_{n=0}^{N-1}\theta_n\right)
 S_{\zeta_0}^{-1}.
 \label{md:comp:telescopaje}
\end{equation}
If chart and phase both return, the result is \(I\). The commutator
in Section~\ref{md:comp:lazo} composes evolutions of different shapes
in a common reference; it is not that passive change.

Under a common representation change \(P\), the operators
\(U,V,H_{\rm lazo}\) are conjugated by \(P\).
The trace and spectrum of \(H_{\rm lazo}\) remain invariant.
If \(X\mapsto PX\) and
\(g\mapsto P^{-\mathsf T}gP^{-1}\) as well, then
\(X^{\mathsf T}gX\) and every correctly transported energy readout are
preserved. Thus \(H_{\rm lazo}\ne I\) cannot be removed by renaming
coordinates.

This constitutes a closed relational mathematical protocol.
Its experimental implementation requires the evolutions and their
inverses to be available physical operations and the system, reference
and controller to have a common realization.
No new force or energy source is asserted here.

\subsubsection{Area conservation and conservation of a common energy}
\label{md:comp:area-energia}

\begin{proposition}[Absence of a common positive metric]
\label{md:comp:prop-metrica}
Each quarter-turn preserves its own metric:
\(g_\zeta=\operatorname{diag}(e^{-\zeta},e^\zeta)\) for \(U_+\),
and \(g_{-\zeta}\) for \(U_-\). When \(\zeta\ne0\), no common
positive metric exists.
\end{proposition}
\begin{proof}
The conditions
\[
 G=\begin{pmatrix}a&b\\b&d\end{pmatrix}>0,
 \qquad U_+^{\mathsf T}GU_+=G
\]
force \(b=0\) and \(d=ae^{2\zeta}\).
The condition \(U_-^{\mathsf T}GU_-=G\) requires
\(d=ae^{-2\zeta}\). The two are compatible only for \(\zeta=0\).
\end{proof}

Preserving symplectic area is therefore not equivalent to preserving
the same energy form throughout every composition of these transports.
In the reference \(g_r=S_r^{-\mathsf T}S_r^{-1}\), with \(\omega_r>0\),
\(X\ne0\) and \(E_r(X)=\omega_r X^{\mathsf T}g_rX/2\), one has
\begin{equation}
 \begin{gathered}
 \frac{E_r(H_{\rm lazo}X)}{E_r(X)}
 =\frac{x^{\mathsf T}\overline{H}_{\rm lazo}^{\mathsf T}
                   \overline{H}_{\rm lazo}x}{x^{\mathsf T}x},\\
 \overline{H}_{\rm lazo}=S_r^{-1}H_{\rm lazo}S_r,\qquad
 x=S_r^{-1}X.
 \end{gathered}
 \label{md:comp:cociente-energia}
\end{equation}
The extrema are the squared singular values.
In the diagonal conjugate case, they are \(e^{\pm8|\zeta|}\),
whereas the state dilations are \(e^{\pm4\zeta}\).
This difference of exponents is necessary. Physically switching
operations requires accounting for the controller; gain in a readout
does not establish energy creation.

\subsubsection{Order requires an oriented readout}
\label{md:comp:orden-orientado}

Interchanging \(U_+\) and \(U_-\) inverts \(H_{\rm lazo}\).
The trace and the set of extremal gains are the same for
\(H_{\rm lazo}\) and \(H_{\rm lazo}^{-1}\). These scalars detect the
loop effect but not its orientation. Two independent inputs with their
complete vector responses reconstruct \(H_{\rm lazo}\); a scalar
measurement is generally insufficient.

In the conjugate case, with marked reference axes, let
\[
 G_Q=\frac{E_r(H_{\rm lazo}e_Q)}{E_r(e_Q)},\qquad
 G_P=\frac{E_r(H_{\rm lazo}e_P)}{E_r(e_P)}.
\]
The readout
\begin{equation}
 \mathcal O=\frac14\log(G_Q/G_P)=4\zeta
 \label{md:comp:lector-orientado}
\end{equation}
recovers orientation and changes sign under loop reversal.
Axes and metric must be transported together under a representation
change. Orientation is thereby distinguished without attributing it
to a trace that removes it.

\subsection{History information versus occupation statistics}
\label{md:comp:informacion}

The products
\[
 U_+U_-U_+^{-1}U_-^{-1}
 \qquad\text{and}\qquad
 U_+U_+^{-1}U_-U_-^{-1}
\]
contain the same four operations with the same multiplicities.
The second equals \(I\); the first is \(H_{\rm lazo}\ne I\) for
\(\zeta\ne0\). Any statistic depending only on multiplicities
therefore assigns the same result to two protocols with different
actions on the state.

The exact conclusion is that such a statistic does not constitute
a sufficient description for predicting the response of the composition.
This does not refute Shannon: a distribution on ordered sequences can
represent order. What is distinguished is the information retained
in the chosen observable.

\begin{definition}[Operational information of a history]
\label{md:comp:def-informacion}
In this representation, the operational information of a history is
proposed to mean the class of its ordered transport distinguishable
by the admitted readouts. Let \(\mathcal D\) be a collection of
detectors defined on operators. It may include
\[
 D_{X,L}(H_{\rm lazo})=L(H_{\rm lazo}X),
 \qquad D_{\rm tr}(H_{\rm lazo})=\operatorname{tr}H_{\rm lazo},
\]
with preparation \(X\) and linear readout operator \(L\).
Two histories are equivalent if
\[
 D(H_{{\rm lazo},1})=D(H_{{\rm lazo},2})
 \qquad\text{for every }D\in\mathcal D.
\]
\end{definition}

Enlarging the collection may separate previously indistinguishable
classes. Complete vector responses on a basis reconstruct every linear
operator; with trace alone,
\(H_{\rm lazo}\sim H_{\rm lazo}^{-1}\) may persist.
This distinguishes the detector of a state response from the spectral
detector of the operator.

The definition does not identify transport with the entire TPK
genealogy: different histories may produce the same operator in a
representation. Information latent in other sheets, incidences and
extensions is retained by its own readout operators.

The statistical and operational dimensions are compatible.
The combinatorial growth rate, or topological entropy, measures
asymptotic growth of admissible configurations; a Shannon entropy
rate additionally requires a distribution.
An observation map determines which transformations can be reconstructed.
None of the preceding proofs asserts universally zero thermodynamic
entropy. The contribution is to specify which description suffices
for a prediction or reconstruction task.

\subsection{Consequences and proposed new definitions}
\label{md:comp:consecuencias}

The preceding results bring together five notions, retaining their
domains and readout operators:
\begin{enumerate}
 \item \textbf{Constitutive spectral potential.}
 \(\mathscr F\), defined by the \(90/120\) moments, has the logarithms
 of normalized speed and normalized impedance as its two derivatives.
 This is a mathematical definition downstream of the channels, with
 no added energy interpretation.

 \item \textbf{Differential constitutive reciprocity.}
 The equality of cross-sensitivities proved in
 Section~\ref{md:comp:reciprocidad} provides a verifiable joint
 constraint in addition to the values.

 \item \textbf{Relative closure defect.}
 \(\operatorname{tr}H_{\rm lazo}-2\) detects the nontrivial loop.
 It is accompanied by an oriented readout, such as \(\mathcal O\),
 when order must be distinguished. Trace alone does not define all memory.

 \item \textbf{Operational reconstruction.}
 This is recovery of the action of a history from a specified
 collection of observables. It is distinguished from two-channel
 reconstruction and from reconstruction of the full state.

 \item \textbf{Reconstruction margin.}
 The bound \(60a(M_0)\) in Section~\ref{md:comp:restitucion}
 quantifies absolute stability on a domain, whereas injectivity
 asserts only exact uniqueness.
\end{enumerate}

The mass, frequency, radius and temperature coordinates of the
preceding development remain unchanged. This development adds two
levels: differential compatibility between scalar readout operators
and order memory in operator compositions.
The dilation factor of \(H_{\rm lazo}\) is not identified with a new
mass without passing through the corresponding mass readout operator.
